% !TEX root = ../main.tex
\section{The Systems in Scope}
\label{sec:architecture}

Before classifying what these systems get wrong, it is worth being exact about
which systems are meant. The argument that follows depends on specific properties
of one architecture and one training regime, and it does not generalise to deep
neural networks at large.

\subsection{What is in scope, and what is not}
\label{sec:scope}

In scope are systems with all four of the following properties.

\begin{enumerate}
  \item They emit natural-language assertions---propositional content a reader can
        act on---as their output.
  \item They produce that output by autoregressive sampling over a fixed token
        vocabulary, one token at a time, each conditioned on the preceding ones.
  \item Their parameters are fitted by next-token likelihood on large text
        corpora.
  \item They are optionally post-trained on human demonstrations and preference
        comparisons.
\end{enumerate}

Concretely: the decoder-only transformer language models established at scale by
the GPT-3 generation \citep{brown2020}, extended by instruction tuning and
preference optimisation \citep{ouyang2022}, and everything built on that pattern
since. That is the point at which these systems became plausible components of a
business process, and it is the regime current deployments inherit.

Out of scope, and worth naming because ``deep learning'' is often treated as one
thing:

\begin{itemize}
  \item \textbf{Generative adversarial networks} and \textbf{diffusion models}.
        Both are generative, and neither generates by autoregressive token
        sampling---an adversarial objective and an iterative denoising objective
        respectively---and neither emits propositional content that an
        organisation relies on as an assertion. A synthesised image makes no claim
        that could be grounded or ungrounded.
  \item \textbf{Encoder-only and discriminative models.} A masked language model
        used as a classifier, a ranker, or an extractor produces a label or a score.
        The exclusion here is one of \emph{tractability}, not of importance, and it
        should not be read as saying such systems raise no governance concern. A
        prediction of ``high risk'', ``eligible'', or ``fraud likely'' is routinely
        acted upon, and it needs provenance, applicability, validation, monitoring, and
        authorization exactly as an assertion does. What differs is that a finite,
        enumerable output space and a task-specific validation regime make the
        technical problem materially different, as Section~\ref{sec:not-easier}
        develops. This paper is about systems whose primary output is open-ended
        natural-language assertion.
  \item \textbf{Conventional predictive and statistical models.} Regression,
        gradient-boosted trees, and the rest of the enterprise analytics estate
        are out of scope for the same reason: they do not assert.
\end{itemize}

\subsection{What the system computes}
\label{sec:what-computed}

The computation has four stages, and the first is where the argument of
Section~\ref{sec:canonical-form} originates.

\paragraph{Tokenisation.} Text is segmented into tokens drawn from a fixed
vocabulary $\mathcal{V}$, typically by byte-pair encoding. This is a purely
lexical operation. It is worth pausing on what it means for the rest of the stack:
\texttt{10}, \texttt{ten}, and \texttt{10.0} are distinct token sequences, as are
\texttt{Q1 2024} and \texttt{first quarter of 2024}. No semantic normalisation
occurs at this boundary, or at any later one.

\paragraph{Embedding.} Token identifiers are mapped to vectors by an embedding
matrix $W_E \in \mathbb{R}^{|\mathcal{V}| \times d}$, and position information is
added. An input of $n$ tokens becomes a matrix in $\mathbb{R}^{n \times d}$.

\paragraph{Causal self-attention.} Each of $L$ layers computes, per head,
\begin{equation}
  \operatorname{Attention}(Q, K, V)
  = \operatorname{softmax}\!\left(\frac{QK^{\top}}{\sqrt{d_k}} + \Lambda\right) V,
  \qquad
  \Lambda_{ij} = \begin{cases} 0 & i \ge j,\\ -\infty & i < j,\end{cases}
  \label{eq:attention}
\end{equation}
where $Q, K, V$ are linear projections of the layer input, $d_k$ scales the logits
to keep the softmax out of saturation, and $\Lambda$ is the causal mask that
prevents position $i$ from attending to positions after it
\citep{vaswani2017}. Attention output is combined with a position-wise
feed-forward transformation, with residual connections and normalisation.

\paragraph{Unembedding and factorisation.} The final layer is projected to logits
over $\mathcal{V}$ and normalised, giving $P(x_k \mid x_{<k}; \Theta)$. The causal
mask is what licenses reading the whole as an autoregressive factorisation:
\begin{equation}
  P(x_1, \ldots, x_n; \Theta) = \prod_{k=1}^{n} P(x_k \mid x_{<k}; \Theta).
  \label{eq:factorisation}
\end{equation}

Now state the signature of the resulting map plainly, because everything in
Section~\ref{sec:canonical-form} follows from it:
\begin{equation}
  M : \mathcal{V}^{*} \longrightarrow \Delta(\mathcal{V}^{*}),
  \label{eq:signature}
\end{equation}
a function from token sequences to distributions over token sequences. The signature is
$M$'s and not $S$'s, and the distinction matters: a deployed system with typed requests,
databases, tools, gates, and escalation paths is a control-flow structure over records and
decisions, not a map from strings to strings. Everything this paper asks for lives in the
difference between the two. The domain
is $\mathcal{V}^{*}$. Not propositions, not normalised queries, not parsed
intents, not records with identity---token sequences. There is no slot in the
parameterisation for a proposition, and no stage between the tokeniser and the
unembedding at which two encodings of the same claim are brought to a common form.
This is not an omission that could be patched at the margins; it is what the
architecture is.

There is direct evidence that content is stored in a form tied to its surface
presentation rather than propositionally. Models trained on statements of the form
``$A$ is $B$'' assert them reliably and fail to answer a query directed at $B$, even
though the fact is the same fact \citep{berglund2024}. The result is about what was
learned rather than about how two phrasings of one request are treated---a
distinction Section~\ref{sec:canonical-form} keeps, and this is not an instance of
the invariance failure discussed there. What it shows is narrower than an
absence of internal representation, and the narrower claim is the one this paper needs.
Whatever representations the model may learn---and Section~\ref{sec:premise} is
deliberately agnostic about them---the architecture exposes no stable, inspectable,
versioned, or contractually enforceable propositional record on which an institution can
rely, and no guaranteed invariance relation between formulations of the same fact. An
assertion's availability depending on the order in which the fact was encountered is
evidence about the second of those, not about what is or is not encoded.

\subsection{The pretraining objective}
\label{sec:pretraining}

Parameters are fitted by minimising cross-entropy on next-token prediction over a
large corpus:
\begin{equation}
  \mathcal{L}_{\mathrm{pre}}
  = -\sum_{k} \log P\!\left(x_k \mid x_{<k}; \Theta\right).
  \label{eq:objective}
\end{equation}
For the GPT-3 generation this was applied at a scale of order $10^{11}$ dense
parameters over a corpus of order $10^{11}$ tokens drawn from filtered web crawl,
curated web text, books, and encyclopaedic sources \citep{brown2020}.

Read \eqref{eq:objective} as a specification and audit it for what is absent. The
claim to make here is narrower than it is tempting to make. Fitting human text does
make the model sensitive to truth-correlated structure---that is why these systems are
useful---so it is wrong to say the objective is blind to truth. What
\eqref{eq:objective} contains is no term that \emph{rewards truth over plausible
falsehood where the two diverge}. On the training distribution they mostly coincide
and the loss cannot tell them apart; off it they separate and the loss remains
indifferent. Beyond that, it contains no term that is a function of evidence. No term that is a function of
evidence, or of the provenance of the text being predicted. No term that is a
function of the context in which an assertion might later be relied upon. And no
term that is minimised by declining to produce a continuation---the objective is
defined over every position, and an abstention is simply a different continuation
with its own probability.

What \eqref{eq:objective} does reward, and rewards very effectively, is assigning
high probability to continuations resembling the training distribution. A fluent
falsehood in a familiar genre resembles the training distribution closely. This is
not a criticism of the objective, which does what it was designed to do; it is an
observation about scope. The system was fitted for \emph{plausible continuation},
and plausible continuation and warranted assertion are different targets that
coincide on the training distribution and separate off it.

One consequence is worth drawing now, because it is the origin of two of the
five failure categories Section~\ref{sec:taxonomy} sets out. Off the training
distribution---an entity never described, a figure never published, a combination
of jurisdiction and effective date that never co-occurred in the corpus---the
system must still emit tokens. It has no representation of the fact that it is now
extrapolating rather than reporting, because \eqref{eq:objective} gave it no
reason to construct one. It produces the most probable continuation, which is
fluent and which rests on nothing. Whether such a continuation happens to be true
is a separate matter, and Section~\ref{sec:taxonomy} argues that the distinction
between those two outcomes is far less important than it looks.

\subsection{In-context learning, and why the prompt is the whole control surface}
\label{sec:icl}

The property of this generation of models that made enterprise use plausible was
not raw perplexity. It was in-context learning: task behaviour induced by
examples and instructions placed in the prompt, with no gradient update
\citep{brown2020}. That is what made a general-purpose model deployable against a
specific business task without a training programme per task, and it is why these
systems entered production at all.

It also has a consequence that is easy to miss and directly relevant to
governance. If task behaviour is induced through the prompt, then the prompt is the
control surface---and by \eqref{eq:signature} the prompt is a token sequence. All
conditioning, including everything an architect would think of as configuration,
policy, or scope, enters the computation as surface form and is processed by the
same mechanism as the content. There is no separate typed channel for constraints.

This is the architectural root of the failure examined in
Section~\ref{sec:canonical-form}. Other channels exist---system prompts,
fine-tuning, adapters, logit biases, tool schemas---but the per-request control
surface, the one that varies with what a user asks, is the prompt. When that surface
is a string, and
the map's domain is strings, sensitivity to the form of the string is not a defect
layered on top of the design. It is the design operating as specified.

\subsection{Instruction tuning and preference optimisation}
\label{sec:rlhf}

Deployed systems add post-training. Supervised fine-tuning on demonstrations is
followed, in the regime established by \citet{ouyang2022}, by preference
optimisation: a reward model $r_\phi$ is fitted to human comparisons between
candidate outputs, and the policy is optimised against $r_\phi$ with a divergence
penalty against the supervised reference,
\begin{equation}
  \max_{\Theta}\ \mathbb{E}\bigl[r_\phi(x, y)\bigr]
  - \beta\, \mathbb{D}_{\mathrm{KL}}\!\left(\pi_\Theta(y \mid x)
    \,\|\, \pi_{\mathrm{ref}}(y \mid x)\right).
  \label{eq:rlhf}
\end{equation}

This changes the output distribution substantially, and it is why current systems
are usable where their pretrained predecessors were not. For the purposes of this
paper, three properties of \eqref{eq:rlhf} matter.

\paragraph{The reward is a function of the output, not of the evidence.}
$r_\phi(x, y)$ is fitted to preferences over $(x, y)$ pairs. Nothing in its domain
is the evidential base, the source version, or the resolved context. An assertion
and its ungrounded twin---identical text, different provenance---are
indistinguishable to the reward model, and therefore receive identical
optimisation pressure.

The precise limitation is worth stating carefully, because the strong version is
false. Preference optimisation \emph{can} select for the observable correlates of
grounding: raters shown a well-cited answer and a bare one will prefer the former,
and instructions to raters can make citation quality part of the comparison. What
$r_\phi$ cannot do is \emph{verify} the relation, because $E$ is not in its domain.
It can learn that citation-shaped output is preferred; it cannot learn that a
citation supports the claim attached to it, except through whatever raters happened
to check at labelling time---which does not generalise to evidence they never saw.
So the pressure is toward the appearance of grounding, and the gap between appearance
and relation is exactly where an auditor has to look.

Two qualifications keep the claim defensible. Deployed systems can place retrieved
evidence, tool output, and structured context inside $x$, and rubric-based or
process-supervised training can score intermediate steps rather than only final
outputs---so the domain of the reward is a design choice rather than a fixed limitation.
And the direction of pressure is an empirical tendency contingent on rater instructions,
reward design, and task incentives, not a theorem: raters told to prefer abstention
where evidence is absent will produce a reward that prefers it. The claim that survives
is conditional. \emph{Preference optimisation alone does not establish a
source-specific, version-specific, context-specific evidence relation for a live
assertion, unless the training and deployment architecture explicitly represents and
verifies that relation.}

\paragraph{The pressure is toward assertion.} Human raters comparing outputs
prefer answers that are fluent, confident, complete, and helpful. A hedged answer,
a request for missing context, or a refusal reads as less helpful than a direct
one, and will usually lose the comparison unless raters are specifically
instructed otherwise. The gradient therefore points away from abstention. The
sycophancy results discussed in Section~\ref{sec:canonical-form} are a measured
consequence of this structure \citep{sharma2024}, not an incidental defect.

\paragraph{The result is better-calibrated form, not better-grounded content.}
Post-training teaches the surface conventions of careful assertion---citation
shapes, hedging vocabulary, structured justifications---because those conventions
are present in preferred outputs. Teaching the relation those conventions signal is
harder and is bounded by what raters could verify: preference data can carry a signal
about whether citations checked out on the examples that were labelled, and cannot
carry one about evidence nobody looked at. So the conventions generalise and the
relation does not. For an auditor this is the least comfortable finding in the section,
because the direction of the gap is unfavourable: post-training reliably improves how
grounded an output \emph{looks} and improves the grounding rate of \eqref{eq:rate}
only incidentally, which widens rather than narrows the distance between apparent and
actual auditability.

\subsection{Three components that are absent}
\label{sec:missing}

Stated as a gap analysis, three things a warranted-assertion system would contain
are absent from the forward pass, and each maps to a category. These are claims about
$M$, not about $S$. A deployed system may well have an evidence gate, a context
resolution step, and an abstention branch---indeed Section~\ref{sec:invariants}
requires all three, as I11, I10, and I17. The point is that none of them comes from the
model, so each must be built, and a system assembled without them inherits the gaps
below by default rather than by choice.

\paragraph{No evidence gate.} Nothing in the computation consults an evidential
base, forms a judgement about whether it supports the candidate claim, and
conditions emission on the answer. Retrieval places evidence \emph{in the
context}, which raises the probability of continuations consistent with it, but
that is influence, not gating: by \eqref{eq:signature} the retrieved text is more
tokens, conditioning a distribution, not a precondition on output. This is why ungrounded
assertion remains available in retrieval-augmented systems, and why the
sensitivity condition of Definition~\ref{def:grounding} is a genuine constraint
rather than a description of existing behaviour.

\paragraph{No context resolution.} Nothing establishes $c$ before answering.
``Is this approved?'' is syntactically complete while leaving jurisdiction,
effective date, material facts, and intended use unspecified, and there is no
mechanism whose function is to notice the omission. This is the failure
Section~\ref{sec:grounding-context} develops as misapplied evidence.

\paragraph{No abstention primitive.} There is no state corresponding to ``I have
not established this.'' Abstention, where it occurs, is a generated
continuation---the tokens of a refusal, produced because refusal-shaped text was
probable in that context, not because a support check failed. This is why
abstention behaviour is sensitive to phrasing: it is an output, not a branch. Any
system that must reliably decline needs the decision to sit outside $M$.

\subsection{A floor, not merely a gap}
\label{sec:floor}

It might be hoped that the ungrounded categories shrink toward zero with better
training. There is a formal reason to expect otherwise.
% TODO(verify): check this gloss against the statement of the theorem before
% submitting. The characterisation below is from the result as I understand it and the
% precise conditions matter.
\citet{kalai2024} show that a calibrated language model must produce incorrect
assertions at a positive rate on facts appearing rarely in its training
distribution: a model well calibrated on the frequency of rare facts is thereby
committed to generating some of them wrongly, and a model that avoided this would
be miscalibrated in a way that costs it elsewhere.

The engineering reading is that the ungrounded cells have a floor above zero that
is not a defect of any particular system, and that the floor is highest exactly
where enterprise questions live---specific entities, specific dates, specific
jurisdictions, all of them rare in a general corpus. If the generation path cannot
be made not to produce ungrounded assertions, then the surrounding system must be
able to \emph{detect} them. That is the argument for measuring grounding directly
rather than waiting for the model to improve.

\subsection{What this predicts about better models}
\label{sec:better-models}

The analysis supports specific predictions, which matter for roadmap decisions
because they are not the predictions usually made.

Scale and post-training reliably reduce outright fabrication on questions
resembling the training distribution. They do not address a corpus that is wrong,
which is a property of the data rather than the model. They do not address
inapplicable evidence with the context it lacks. Capability gains can make a model
more likely to \emph{ask} for a missing jurisdiction or effective date, which helps;
they cannot make asking reliable, and they cannot supply the answer, which has to come
from the request or from a system that resolves it. And their effect on lucky
guessing runs in the wrong direction: a better model guesses better, which raises
accuracy while leaving the assertion exactly as disconnected from evidence as
before. Section~\ref{sec:taxonomy} gives these four outcomes names and
Section~\ref{sec:remedies} assigns them owners.

One concession is owed here, and it is specific. A system that samples several
candidate outputs and selects among them with a verifier does have a check outside the
forward pass, and where the verifier is sound---a proof checker, a test suite, a
constraint solver---that check is a genuine correctness oracle. Such a system satisfies
something structurally similar to the external gate this paper argues for in
Section~\ref{sec:constrained-output}, and it would be wrong to claim the mechanism is
absent from systems that plainly have it.

What does not transfer is the availability of the oracle. Verification of that kind
requires correctness to be cheaply decidable against a formal specification, which is
why it works for mathematics, code, and constraint satisfaction. There is no oracle
for whether a particular regulatory instrument is the governing one for this
jurisdiction at this effective date, applied to this population, for this decision
class. The question is not undecidable in some deep sense; it is simply not settled by
computation over the assertion. It is settled by evidence and by authority, which is
what Sections~\ref{sec:grounding-context} and~\ref{sec:governed-rag} are about. So
outcome verification is a real mechanism that addresses a class of problems this paper
is not about, and its success there is not evidence about the class it is.

Longer context windows, retrieval, and tool use change what the system can do without
changing \eqref{eq:signature}. The domain
remains $\mathcal{V}^{*}$; there is still no canonicalisation stage; the reward
still cannot see provenance. Retrieval is the most consequential of them for this
paper's purposes, and its contribution is not primarily accuracy---it is that an
identified evidential base makes the grounding rate observable at all, as
Section~\ref{sec:measurability} explains.

\subsection{How this compares with other model families}
\label{sec:not-easier}

Section~\ref{sec:scope} excluded discriminative and predictive models on the ground
that grounding is not a coherent question about a label. With the architecture now in
view the comparison can be made properly, and it should be, because the exclusion is
about \emph{whether grounding applies} and not about whether those systems are easy to
validate. It would be wrong to suggest that discriminative
and predictive models are unproblematic, and the overlap with this paper's categories
is worth naming, because it locates the argument honestly.

A classifier that has learned a spurious feature---the hospital identifier rather
than the pathology, the formatting artefact rather than the content---is right for
the wrong reason. That is the same defect this paper calls a \emph{lucky} assertion:
correct output, no supporting relation, and no reason to expect the next case to go
the same way. A classifier whose label flips under an imperceptible perturbation has
failed a form-invariance requirement. A classifier degrading under distribution
shift has exactly the problem of a sample that does not generalise. And per-instance
attribution is weak for most model families, which is a traceability failure.
Enterprise ML monitoring exists because none of this is solved.

What is different is narrower and it is about whether the validation problem is
\emph{well posed}. For a classifier it is. The output space is finite and
enumerable, so a complete reference set can in principle be constructed. Errors are
label mismatches: atomic, countable, and unambiguous. The input space carries a
metric, so perturbation robustness can be stated as a radius and, for some model
classes, certified within it. Residual risk concentrates in distribution shift,
which is a monitorable quantity with established practice around it.

None of those four properties holds for an assertion. The output space is
$\mathcal{V}^{*}$, so there is no complete reference set to build. An error is not a
label mismatch: an assertion can be right in one jurisdiction and wrong in another,
adequate for a weaker claim than the one made, or true for reasons unconnected to
its evidence. The equivalence classes on the input are semantic rather than metric,
so there is no radius to certify and no coverage strategy against a class you cannot
enumerate. And the assertion makes a claim, so its correctness is relative to a
context that the input may never have specified.
